\section{相关工作}
\label{sec:related}

本文的工作处于三条研究与工程脉络的交汇处：AI 编码代理、终端多路复用与自动化、以及 LLM 工具调用与自主代理循环。我们逐一梳理并指出 WhatyTerm 的差异化定位。

\subsection{AI 编码代理}

以 SWE-agent~\cite{yang2024sweagent}、Devin 类端到端软件工程代理为代表的一批工作，聚焦于让 LLM 自主完成从问题理解到代码修改的闭环，并在 SWE-bench~\cite{jimenez2024swebench} 等基准上评估其解决真实 GitHub issue 的能力。这类系统通常自建“代理—环境”接口（Agent-Computer Interface），把文件编辑、命令执行等抽象为代理可调用的动作空间。与之并行，工业界的命令行编码代理（Claude Code、Codex CLI、Gemini CLI 等）则直接以人类使用的终端为运行环境。

WhatyTerm 与这两类工作的根本区别在于\textbf{关注层次}。SWE-agent 类工作关注“代理如何决策与修改代码”，而 WhatyTerm 不重新实现一个编码代理，而是关注“如何让\emph{已有的}命令行编码代理在无人值守下持续运行”。换言之，WhatyTerm 是编码代理之上的\emph{编排与监督层}：它把多个异构 CLI（Claude/Codex/Gemini 等）视作黑盒执行体，通过感知其终端状态并注入操作来消除人工中断点。这一定位使其与具体编码代理解耦，可同时驱动多家 CLI。

\subsection{终端多路复用与自动化}

tmux、GNU screen 等终端多路复用器提供了会话持久化、断线重连与程序化控制（如 \texttt{send-keys}、\texttt{capture-pane}）能力，长期用于服务器运维与自动化脚本。传统的终端自动化工具（如 expect）则以“模式匹配—发送响应”的方式驱动交互式程序。

WhatyTerm 复用了 tmux 作为会话持久层——真正的代理进程活在 tmux 会话中，前端 PTY 仅作为可观察窗口，从而实现长任务的断连续跑与可回看。但其状态判断不再是 expect 式的静态模式匹配，而是\textbf{规则状态机与 LLM 语义分析相结合}的动态决策，且面对的是 LLM 编码代理特有的、由富文本 TUI 渲染的复杂屏幕。这使得 WhatyTerm 在“程序化终端控制”这一传统能力之上，叠加了针对 AI 代理交互模式的语义理解层。

\subsection{LLM 工具调用与自主代理循环}

ReAct~\cite{yao2023react} 提出的“推理—行动”交替范式，以及围绕工具调用（function calling）构建的各类自主代理框架，奠定了 LLM 驱动多步任务的基础。此类“感知—决策—行动”循环是自主代理的通用骨架，但朴素实现往往缺乏对\emph{长程任务全局结构}的组织，容易陷入重复动作或无界循环。

WhatyTerm 在两个方面深化了这一范式。其一，在\textbf{决策层}引入分层结构，以确定性规则消化高频确定性状态，仅在不确定时调用 LLM，兼顾成本与可控性——这与“每一步都调用 LLM”的朴素循环形成对比。其二，在\textbf{任务层}引入三层嵌套循环，通过需求分解、接口契约、拓扑排序并行、集成验证与失败根因分类，为长程任务施加显式的全局结构，并以每层独立的重试预算保证终止性。据我们所知，将“失败类型的语义分类”用于选择“最小修复范围（局部重跑 vs.\ 回退重拆）”，是本系统区别于通用代理循环的一个具体设计。

\subsection{小结}

综上，WhatyTerm 的差异化定位可概括为：它不是又一个编码代理，而是一个\textbf{面向异构命令行编码代理的、以分层决策与多层自主循环为核心的无人值守编排系统}，在传统终端多路复用能力之上叠加了 AI 语义感知，在通用代理循环之上叠加了成本可控的分层决策与结构化的长程任务组织。
